% --- Archon physics formalization source begin ---
% archon:physics
% archon:chemistry
% archon:covers IChO2026Problems/problem_icho_2026_t3_a2.lean
% archon:source-report reports/icho_2026/problem_icho_2026_t3_a2.source.json
% archon:problem-id icho_2026_t3
% archon:part-id T3-A2

\chapter{Icho Chemistry problem icho\_2026\_t3\_a2}
\label{ch:IChO2026Problems_problem_icho_2026_t3_a2}

\paragraph{Problem source.}
\#\# Source
58th International Chemistry Olympiad, Tashkent, Uzbekistan, 2026.
Official-materials index: https://www.icho2026.uz/problems
English problem PDF: ../icho\_2026\_source/raw/theory\_problem.pdf
English solution/rubric PDF: ../icho\_2026\_source/raw/theory\_solution.pdf
Problem PDF page 26; printed local question page 2; solution starts on PDF page 23.
The page PNG is primary visual evidence for formulas, structures, tables, and figures.

\#\# T3. Into Reticular Chemistry
T3. Into Reticular Chemistry 7\% Covalent organic frameworks (COFs) are a class of polymers that form two- or three-dimensional structures through reactions between organic precursors, resulting in covalent bonds to afford porous, crystalline materials. Two examples of honeycomb type 2D-COFs formed by boronate ester and boroxine rings (repeat units indicated using dashed lines) are shown in the figure below. The black and red bold lines represent different building blocks in all COF topology figures. 3.1 Determine the empirical formula of COF-1 and calculate the mass percentage (\%, to two decimal places) of carbon in COF-1.

\#\# Current subquestion 3.2
Calculate the internal diameter, d, in Å, of the honeycomb for COF-2, with the following assumptions: • C–C/C=C (arenes) = 1.39 Å • C–B = 1.56 Å • B–O = 1.38 Å • Note, the diameter of a circle, 𝑑, inscribed inside a hexagon, is given by 𝑑= √3𝑎, where 𝑎is the side length of the hexagon. • Neglect the width of the linkers.

\#\# Dataset metadata
IChO 2026; T3; raw rubric points=3; kind=theory; formalization\_ready=true.

\paragraph{Shared problem context.}
T3. Into Reticular Chemistry 7\% Covalent organic frameworks (COFs) are a class of polymers that form two- or three-dimensional structures through reactions between organic precursors, resulting in covalent bonds to afford porous, crystalline materials. Two examples of honeycomb type 2D-COFs formed by boronate ester and boroxine rings (repeat units indicated using dashed lines) are shown in the figure below. The black and red bold lines represent different building blocks in all COF topology figures. 3.1 Determine the empirical formula of COF-1 and calculate the mass percentage (\%, to two decimal places) of carbon in COF-1.

\paragraph{Current subquestion.}
Calculate the internal diameter, d, in Å, of the honeycomb for COF-2, with the following assumptions: • C–C/C=C (arenes) = 1.39 Å • C–B = 1.56 Å • B–O = 1.38 Å • Note, the diameter of a circle, 𝑑, inscribed inside a hexagon, is given by 𝑑= √3𝑎, where 𝑎is the side length of the hexagon. • Neglect the width of the linkers.

\paragraph{Recorded answer/context.}
For COF-2, 𝑎= 2 ⋅𝐵−𝑂+ 2 ⋅𝐶−𝐶+ 2 ⋅𝐶−𝐵= 8.66 𝑑= √3𝑎= 15.0 𝑑= 2𝑟= 15.0 d = 15.0 Å 3 pt for correct d 2 pt if only a is given 1 pt if only correct formula is given 3.0 pt

\paragraph{Figure/image paths.}
\begin{itemize}
\item ../icho\_2026\_source/image/T3\_page-2.png
\item ../icho\_2026\_source/image/T3\_page-1.png
\end{itemize}

\paragraph{Verified structural interpretation.}
COF-2 is formed by A2 self-condensation into boroxine rings.  Along one pore
edge the path O--B, B--C, C--C, C--C, C--B, B--O gives two bonds of each
supplied type, hence $a=8.66$ Å and $d=\sqrt{3}\,a$, within $0.05$ Å of
$15.0$ Å.

\paragraph{Formalization target.}
The verified Lean formalization is in `IChO2026Problems/problem\_icho\_2026\_t3\_a2.lean`,
with its theorem and lemma proofs complete.
The Lean declarations must preserve every mathematical object, quantifier, hypothesis, side condition, approximation/error guarantee, and requested conclusion from the source statement.
The implementation reuses the pinned Mathlib, Physlib, and Chemistry libraries,
with local definitions only for source-specific objects.

\begin{theorem}[Icho Chemistry formalization target]
\lean{IChO2026.T3.A2.internal_diameter_cof2}
\leanok
\label{thm:physics:icho_2026_t3_a2:target}
The mapped declarations encode the stated calculation and its verified conclusion.
\end{theorem}
\begin{proof}

\leanok
The proof unfolds the source-specific definitions and establishes the mapped
conclusion from the encoded hypotheses.
\end{proof}
% --- Archon physics formalization source end ---
